% chapters/design/api_design.tex
% Sección 4.3: Diseño de la API REST (~1-2 páginas)
% ============================================================================
%
% GUÍA: Describir los principios de diseño de la API y su estructura.
%
%   1. Principios de diseño:
%      - RESTful: recursos, verbos HTTP, códigos de estado estándar.
%      - Versionado de API (si aplica): /api/v1/...
%      - Paginación offset-based (\ref{rnf:perf:pagination}).
%      - Filtrado y ordenación.
%      - Códigos de error estandarizados (\ref{rnf:maint:errors}).
%      - Serialización con DRF: serializers, viewsets, routers.
%
%   2. Estructura de endpoints (tabla resumen):
%      Agrupar por recurso. No listar todos, solo los más representativos.
%
%      \begin{table}{tab:api_endpoints}{Resumen de endpoints principales de la \ac{api}}
%        \begin{tabular}{l l l}
%          \textbf{Recurso}  & \textbf{Método} & \textbf{Descripción} \\
%          \hline
%          /auth/login        & POST & Inicio de sesión \\
%          /auth/logout       & POST & Cierre de sesión \\
%          /users/            & CRUD & Gestión de usuarios \\
%          /subjects/         & CRUD & Gestión de asignaturas \\
%          /exams/            & CRUD & Gestión de exámenes \\
%          /exam-instances/   & GET/PATCH & Instancias de examen \\
%          /annotations/      & CRUD & Anotaciones de corrección \\
%          ...
%        \end{tabular}
%      \end{table}
%
%   3. Documentación automática:
%      - drf-spectacular para generación OpenAPI 3.0.
%      - Interfaz Swagger UI y Redoc accesibles (\ref{rnf:test:openapi}).
%
%   4. Seguridad de la API:
%      - Rate limiting (\ref{rnf:sec:ratelimit}).
%      - CORS (\ref{rnf:sec:cors}).
%      - Validación de entradas (\ref{rnf:sec:input}).
%
%   NOTA: Los detalles completos de cada endpoint se pueden consultar
%   en la especificación OpenAPI generada automáticamente.
% ============================================================================

% TODO: Redactar la sección del diseño de la API.
